% !TeX spellcheck = pt_PT

\chapter{Declaração de Utilização de Inteligência Artificial}
\label{chap:declaracao-ia}

A utilização de inteligência artificial (IA) constituiu um eixo transversal do presente trabalho, alinhado com uma das propostas orientadoras do projeto: explorar formas de acelerar o desenvolvimento em colaboração com ferramentas de IA, mantendo simultaneamente a qualidade, a integridade e a responsabilidade técnica sobre o código e o texto produzidos.

Ao longo de todas as fases do projeto — desde a conceção inicial e definição de arquitetura até à implementação, testes, documentação técnica, configuração de pipelines de integração contínua e redação deste relatório — recorreu-se a assistentes baseados em modelos de linguagem de grande dimensão (\textit{Large Language Models}). As ferramentas utilizadas incluem o GitHub Copilot, o ChatGPT, o Codex e o Claude.

Estas ferramentas foram empregues, nomeadamente, para sugestão e geração de código, apoio à escrita de testes automatizados, explicação de erros e \textit{debugging}, elaboração e reformulação de texto técnico, tradução de conteúdos, revisão de lógica de negócio e aceleração de tarefas repetitivas ou de maior volume documental. Em todos os casos, os resultados obtidos foram revistos, validados e adaptados por nós ( autores ), e assim, assumimos a responsabilidade final pelo conteúdo integrado no sistema e descrito neste documento.

A abordagem adotada não substituiu o raciocínio crítico nem as decisões de engenharia: tratou-se de um trabalho em conjunto com a IA, em que mantivemos o controlo sobre as escolhas arquiteturais, a conformidade com os requisitos funcionais, a execução de testes e a verificação do comportamento do sistema em ambiente real. O objetivo foi aumentar a produtividade sem comprometer a fiabilidade da solução desenvolvida.
